%% ma: L12-B02-0009
%% lop: 12
%% chuong: 1
%% bai: 2
%% dang: 12.2.6
%% loai: TN4
%% muc: TH
%% dap_an: "D"
%% nguon: "Ảnh giáo viên gửi 11/10/2026 (Đề kiểm tra giữa HK1 Toán 12, THPT Việt Đức 2026-2027, mã đề 0103), Phần I Câu 12"
%% kien_thuc: "Bài 2 (giá trị lớn nhất, nhỏ nhất trên một đoạn đọc từ đồ thị)"
%% trang_thai: chinh-thuc
%% ghi_chu: "Đã duyệt 12/10/2026. Hình TikZ vẽ f(x) = 3 - (x-1)^2 trên [0;3] (đỉnh (1;3), qua (0;2) và (3;-1)) rồi đoạn thẳng từ (3;-1) đến (4;4), theo các số liệu đọc từ hình gốc."
\begin{ex}
Cho hàm số $y=f(x)$ liên tục trên $[0;4]$ và có đồ thị như hình vẽ.
\begin{center}
\begin{tikzpicture}[x=1.4cm,y=.8cm]
  \draw[truc] (-.4,0)--(4.7,0) node[below]{$x$};
  \draw[truc] (0,-1.6)--(0,4.8) node[right]{$y$};
  \draw[thick,domain=0:3,samples=60] plot (\x,{3-(\x-1)*(\x-1)});
  \draw[thick] (3,-1)--(4,4);
  \draw[phu] (1,0)--(1,3)--(0,3) (3,0)--(3,-1)--(0,-1) (4,0)--(4,4)--(0,4);
  \foreach \t in {1,2,3,4} \draw (\t,.08)--(\t,-.08) node[below,font=\footnotesize]{$\t$};
  \foreach \t in {-1,1,2,3,4} \draw (.05,\t)--(-.05,\t) node[left,font=\footnotesize]{$\t$};
  \node[below left,font=\footnotesize] at (0,0) {$O$};
  \node[font=\footnotesize] at (2.45,1.2) {$y=f(x)$};
\end{tikzpicture}
\end{center}
Tổng của giá trị lớn nhất và giá trị nhỏ nhất của hàm số trên đoạn $[0;3]$ bằng
\choice{$3$}{$4$}{$1$}{\True $2$}
\loigiai{Trên đoạn $[0;3]$, điểm cao nhất của đồ thị có tung độ $3$ (tại $x=1$) và điểm thấp nhất có tung độ $-1$ (tại $x=3$). Do đó $\max_{[0;3]}f(x)=3$, $\min_{[0;3]}f(x)=-1$ và tổng của chúng bằng $2$.}
\end{ex}
